\documentclass{article}
\usepackage{pathfinder-paper}

\title{A Reciprocal Contract Witness for Asymmetric CT-Bench Boards}
\pathfinderpapers{2609.01595}{Mechanism Design for Alignment and Control}%
  {2607.22750}{Commitment To Cooperation With Self-Negotiated Contracts}

\begin{document}
\maketitle

\begin{abstract}
We give a finite, computationally checked contract-and-path witness for the asymmetric board class of
CT-Bench. For each of its \(396\) boards, a five-clause programmatic tile contract and a shared shortest route
make both players strictly exceed their non-interaction scores while attaining the game's maximum joint score.
An accounting identity states the delivered net coverage Red needs to improve on its outside option. These are
feasibility statements. They do not establish that agents will propose or accept the contract, or that its
path is an equilibrium of the sequential game. % ledger: 20, 22, 24, 26
\end{abstract}

\section{Question and setting}

Q distinguishes feasible action rules from truthful reporting and obedience, but its one-shot formulation does
not give an equilibrium theorem for P's sequential CT-Bench game \cite{bergemann2026mechanism}. P measures how
language-model agents negotiate and use contracts there \cite{wyse2026commitment}. We ask whether its
asymmetric boards admit a strictly mutually beneficial outcome through its tile-contract interface.
% ledger: 1, 2, 18, 25

Each player starts in one corner of P's \(4\times4\) board and seeks the opposite green goal. The other
\(14\) cells split equally between Red and Blue. Red starts with \(14\) Red and two green chips, Blue with
\(14\) Blue and two green chips. A move consumes one chip matching its destination; a finisher earns \(20\)
plus \(5\) per remaining chip, and a non-finisher earns zero. In the asymmetric class only Red has an
own-chip route. Outside scores are \((b_R,b_B)=(70,0)\), maximum joint score is \(140\), and P samples \(40\)
boards from this class \cite{wyse2026commitment}. % ledger: 11, 20, 22, 24

P's programmatic tile contract (PTC) specifies a tile, giver, receiver and colour per clause. After acceptance,
a covered move automatically transfers a chip if the giver has one; otherwise the giver scores zero. The
construction uses this format and sufficient inventory \cite{wyse2026commitment}. % ledger: 20, 21

\section{Finite construction}

The recorded exhaustive check enumerates all \(\binom{14}{7}=3{,}432\) balanced boards. Exactly \(396\) have
independent Red and dependent Blue. A shortest Blue path has six moves: five onto coloured cells and one onto
the green goal. Among shortest Blue paths, the minimum number \(k\) of Red cells is one on \(314\) boards and
two on \(82\). A separate enumeration of all simple paths found no longer path with fewer Red cells. These
counts are finite computational findings, rather than a closed-form count. % ledger: 11, 14, 20

Here is the reproducible board test. Number cells rowwise from \(0\) (start) to \(15\) (goal), and let
\(R\subseteq\{1,\ldots,14\}\) be the seven Red cells; all other nonterminal cells are Blue. Let \(\mathcal P\)
contain every simple, orthogonally adjacent path from \(0\) to \(15\), and let \(I(p)\) exclude its endpoints.
The asymmetric predicate is
\[
 \bigl(\exists p\in\mathcal P: I(p)\subseteq R\bigr)
 \quad\text{and}\quad
 \bigl(\nexists p\in\mathcal P: I(p)\cap R=\varnothing\bigr).
\]
This is P's Red-independent, Blue-dependent class under its fixed initial chips. % ledger: 11, 14, 20

The enumeration proceeds as follows. % ledger: 11, 14, 20, 21
\begin{enumerate}
\item Generate \(\mathcal P\) by depth-first search without revisiting a cell, then retain its six-move paths as \(\mathcal P_6\).
\item For each of the \(\binom{14}{7}\) choices of \(R\), apply the predicate above.
\item For each retained board compute \(k=\min_{p\in\mathcal P_6}|I(p)\cap R|\), and check
\(k=\min_{p\in\mathcal P}|I(p)\cap R|\); count boards by \(k\).
\item For a shortest path attaining \(k\), check the reciprocal contract inventories and scores below.
\end{enumerate}
The runnable supplementary file \texttt{asymmetric\_oracle.py}, distributed with this paper, implements these
steps and prints the counts and score pairs. % ledger: 11, 14, 20, 21

\begin{proposition}[Reciprocal tile-contract witness]
For each of the \(396\) asymmetric boards under P's stated generation rule, there is a PTC with five
coloured-tile clauses and a pair of six-move paths whose executed scores are \((85,55)\) or \((75,65)\). Both
players strictly exceed \((70,0)\), and their joint score is \(140\). % ledger: 20, 21, 26
\end{proposition}

\begin{proof}
Choose a shortest Blue path minimising its Red-cell count \(k\in\{1,2\}\). Have both players traverse that
same path. For each of its \(k\) Red cells, the contract makes Red cover Blue's move; for each of its \(5-k\)
Blue cells, it makes Blue cover Red's move. Each player pays for its own matching-colour moves and its own
green goal move. Red therefore spends \(2k+1\) chips and Blue spends \(2(5-k)+1\). Neither inventory is
exhausted. As both begin with \(16\) chips, their scores are
\[
(r_R,r_B)=\bigl(20+5[16-(2k+1)],\;20+5[16-(2(5-k)+1)]\bigr)
          =(95-10k,\;45+10k).
\]
The checked values of \(k\) yield the stated pairs. Both paths have the minimum six moves, so the sum is P's
joint maximum of \(140\). % ledger: 20, 21, 26
\end{proof}

This is an existence result conditional on the contract being accepted and its specified route being followed.
Automatic transfer removes an execution choice at each contracted tile when inventory is sufficient. It does
not establish incentives to negotiate these clauses or full sequential obedience after acceptance.
% ledger: 20, 21, 25, 26

\section{Delivered coverage and observed outcomes}

\begin{proposition}[Red's score identity]
Suppose Red finishes in \(L_R\geq6\) moves, Blue pays for \(c\) of Red's moves, and Red pays for \(d\) of
Blue's. With no other chip trades or point transfers,
\[
r_R=70+5\bigl[c-d-(L_R-6)\bigr].
\]
Thus Red strictly beats its own outside score precisely when \(c-d>L_R-6\). % ledger: 22, 23, 26
\end{proposition}

\begin{proof}
Red's remaining inventory is \(16-L_R+c-d\). Applying P's terminal score rule gives \(20+5(16-L_R+c-d)\),
which rearranges to the identity. The proposition concerns Red's score; both-beat-baseline additionally
requires Blue to finish with a positive score. For the reciprocal witness, \(c=5-k\), \(d=k\) and \(L_R=6\),
so the inequality holds for both recorded values of \(k\). % ledger: 22, 23, 26
\end{proof}

P already counts \emph{net tiles promised} in PTC agreements and finds equal numbers assigned to each player
in \(49\%\) of asymmetric contracts. Its PTC visibility intervention moves Red's mean net promised tiles from
\(0.2\) to \(0.45\), without a statistically decisive difference in the reported test. The identity adds a
route-length adjustment and requires \emph{delivered} coverage: promised clauses may not all be triggered, and
later coverage may change the net flow. Accordingly, \(49\%\) is not an estimate of failures caused by
balanced terms \cite{wyse2026commitment}. % ledger: 22, 23

P reports mean both-beat-baseline of \(0.06\pm0.03\) on asymmetric boards for Pay-for-Partner with PTC. Its
\(0.80\pm0.03\) contract-acceptance figure is the all-board mean, across independent, mutually dependent and
asymmetric boards; the two figures therefore have different board denominators. These are aggregate results
of its sampled evaluations. The proposition gives a
feasible score on each board in the generation class, but the contrast cannot identify whether agents failed
to find suitable terms, accept them, traverse the route, or realise the transfers \cite{wyse2026commitment}.
% ledger: 24, 25, 26

\section{Related work and interpretation}

P itself defines the baseline and joint-score metrics, describes the tile-contract mechanism, and analyses
promised net tiles. Earlier Colored Trails work also uses a computed Pareto frontier as an evaluation
benchmark in a different three-player game \cite{wyse2026commitment,li2025combining}. Thus a Pareto benchmark
and the value of net promises are prior ingredients. The contribution here is the recorded all-board witness
for P's particular asymmetric class and its delivered-coverage score test. The search found no published
statement of that exact construction; this is a scoped literature observation, not a proof of priority.
% ledger: 20, 22, 23, 25; related-work search: search.md

Q's feasibility versus incentive distinction explains the interpretation: the oracle removes the possibility
that these boards lack a mutually improving contract-and-path outcome, while leaving the behavioural cause of
P's low observed rate unresolved \cite{bergemann2026mechanism,wyse2026commitment}. The local points-contract
calculation in the peer record gives \(x(p_1-p_0)\geq5\) for honouring a costly coverage at a fixed history;
it is a continuation comparison, not a sequential equilibrium condition. In particular, the \(20\)-point cap
alone does not rule out the one- or two-coverage oracle paths. % ledger: 9, 10, 14, 15, 18, 25

\section{Limitations and open questions}

The enumeration covers P's stated balanced \(4\times4\) generation class, including its asymmetric rule. It is
a finite computational check, not a theorem for larger boards or changed chip endowments. The construction
assumes the compiled tile clauses are accepted, enough chips remain when invoked, and both players use the
selected path. It proves neither a bargaining equilibrium nor sequential incentive compatibility under
arbitrary deviations. P's aggregate table does not disclose which stage blocked the oracle outcome on each
run. % ledger: 11, 14, 20, 21, 24, 25

A discriminating experiment would supply the verified contract as already accepted on each asymmetric board,
allow autonomous movement, and log route lengths and delivered transfers. Comparing that condition with
autonomous proposal and acceptance would separate finding the terms from executing them; replaying the same
coverage as defeasible Pay-for-Partner would test the role of automatic enforcement. These are proposed tests,
not observed effects. Even success in the supplied-contract condition would leave voluntary bargaining and
full sequential incentives open. % ledger: 18, 25, 26

\bibliographystyle{plainurl}
\bibliography{references}
\end{document}
